\documentclass[conference]{IEEEtran}
\IEEEoverridecommandlockouts
% The preceding line is only needed to identify funding in the first footnote. If that is unneeded, please comment it out.
\usepackage{cite}
\usepackage{amsmath,amssymb,amsfonts}
\usepackage{algorithmic}
\usepackage{algorithm}
\usepackage{graphicx}
\usepackage{textcomp}
\usepackage{float}
\usepackage{xcolor}
\def\BibTeX{{\rm B\kern-.05em{\sc i\kern-.025em b}\kern-.08em
    T\kern-.1667em\lower.7ex\hbox{E}\kern-.125emX}}
\begin{document}

\title{AoI-Aware Task Offloading for MEC Using TLBO-HHO Algorithm}

\author{\IEEEauthorblockN{1\textsuperscript{st} Given Name Surname}
\IEEEauthorblockA{\textit{dept. name of organization (of Aff.)} \\
\textit{name of organization (of Aff.)}\\
City, Country \\
email address or ORCID}
\and
\IEEEauthorblockN{2\textsuperscript{nd} Given Name Surname}
\IEEEauthorblockA{\textit{dept. name of organization (of Aff.)} \\
\textit{name of organization (of Aff.)}\\
City, Country \\
email address or ORCID}
\and
\IEEEauthorblockN{3\textsuperscript{rd} Given Name Surname}
\IEEEauthorblockA{\textit{dept. name of organization (of Aff.)} \\
\textit{name of organization (of Aff.)}\\
City, Country \\
email address or ORCID}
}

\maketitle

\begin{abstract}
\textbf{Mobile Edge Computing (MEC) task offloading faces the challenge of optimizing energy consumption, latency, and information freshness simultaneously. This paper proposes a two-phase integrated TLBO-HHO algorithm that combines Teaching-Learning-Based Optimization and Harris Hawks Optimization to address this tri-objective optimization problem. We establish an optimization model incorporating Age of Information (AoI) as a key metric for information freshness, prove its NP-hardness, and develop an adaptive switching mechanism that balances global exploration and local exploitation. Experimental results demonstrate that TLBO-HHO achieves 65.2\%, 85.4\%, and 74.6\% improvement in fitness values compared to TLBO, GA, and GWO algorithms respectively, while maintaining 95\% of real-time tasks' AoI within 1 second through intelligent resource allocation strategies.}
\end{abstract}

\begin{IEEEkeywords}
Mobile edge computing, task offloading, resource allocation, age of information, TLBO-HHO algorithm.
\end{IEEEkeywords}

\section{Introduction}
\IEEEPARstart{W}{ith} the rapid development of Internet of Things (IoT) and 5G/6G communication technologies, the amount of data generated by mobile devices is growing exponentially. Mobile Edge Computing (MEC) provides an effective solution for delay-sensitive applications by deploying computing resources at the network edge~\cite{b1}. In MEC systems, task offloading decisions directly affect system performance and need to balance between energy consumption and latency. In recent years, with the increasing demand for real-time communication and real-time decision-making, Age of Information (AoI) has become an important metric for measuring system information freshness~\cite{b15}, and incorporating AoI optimization has become a research trend.

Existing research methods include convex optimization~\cite{b2}, game theory, deep reinforcement learning~\cite{b3}, and metaheuristic algorithms. However, single metaheuristic algorithms have inherent defects: Genetic Algorithm (GA) is prone to premature convergence, and Grey Wolf Optimizer (GWO) has slow convergence speed~\cite{b4}. Hybrid metaheuristic algorithms can effectively overcome these limitations by combining the advantages of different algorithms.

The main contributions of this paper are as follows:
\begin{itemize}
    \item We establish a tri-objective optimization model considering energy consumption, latency, and AoI, with weight settings of 0.25-0.35-0.40, analyze task waiting time through queuing theory, and prove its NP-hard nature.
    \item We propose a two-phase integrated TLBO-HHO fusion algorithm that achieves dynamic balance between exploration and exploitation through an adaptive probability mechanism (0.7$\rightarrow$0.1).
    \item We validate the effectiveness of the algorithm in information timeliness optimization in a scenario containing 20 devices, 4 edge servers, 2 cloud servers, and 40 tasks.
\end{itemize}

\section{Related Works}

The task offloading optimization problem in MEC has attracted significant research attention, with various approaches ranging from mathematical optimization to artificial intelligence techniques. Among these, metaheuristic algorithms have emerged as promising solutions due to their ability to handle non-convex optimization problems effectively.

Teaching-Learning-Based Optimization (TLBO), proposed by Rao~\cite{b5}, simulates the classroom teaching process and has demonstrated advantages of parameter-free operation and stable convergence. Similarly, Harris Hawks Optimization (HHO), introduced by Heidari et al.~\cite{b6}, mimics the cooperative hunting behavior of Harris hawks and achieves success rates of 85-95\% when dealing with multimodal functions. Recent advances have seen the hybridization of these algorithms; Chen et al.~\cite{b7} and Tu et al.~\cite{b8} developed improved versions of HHO by combining it with other metaheuristics. However, the specific combination of TLBO and HHO for MEC applications remains unexplored, presenting an opportunity for innovation.

In the context of MEC task offloading, existing optimization methods can be broadly categorized into three approaches. Mathematical optimization methods provide theoretical guarantees but struggle with practical constraints. Wang et al.~\cite{b9} developed a joint computation offloading scheme based on Lagrangian dual decomposition, while Liu et al.~\cite{b10} proposed a Lyapunov optimization-based algorithm for delay-optimal scheduling. Although these methods offer optimal solutions under certain conditions, they often fail to handle the non-convex nature of real-world MEC systems. Machine learning approaches, particularly deep reinforcement learning, have shown promise in adaptive resource allocation. Wang et al.~\cite{b11} demonstrated dynamic trajectory control in UAV-assisted MEC, and Zang et al.~\cite{b13} integrated federated learning for privacy-preserving optimization. Nevertheless, these methods require extensive training data and computational resources. Metaheuristic algorithms bridge this gap by offering practical solutions without extensive training requirements. Li and Zhu~\cite{b12} achieved 30\% energy reduction using genetic algorithms, while Ali et al.~\cite{b14} reported up to 56\% energy savings with multi-objective HHO in cloud-fog computing scenarios.

A critical yet underexplored dimension in MEC optimization is the Age of Information (AoI), which quantifies information freshness at the receiver. First introduced by Kaul et al.~\cite{b16} in 2012, AoI has gained increasing relevance in IoT and real-time applications. Abd-Elmagid et al.~\cite{b15} emphasized its importance in status update systems, highlighting the need for AoI-aware optimization in time-sensitive applications. Recent studies have begun incorporating AoI into MEC frameworks; Liu et al.~\cite{b17} developed an asynchronous deep reinforcement learning method for AoI-aware task computing in vehicular networks, while Zhou et al.~\cite{b18} designed contract matching mechanisms that consider information freshness. Despite these advances, the integration of AoI optimization with metaheuristic algorithms for MEC task offloading remains largely unexplored. This gap motivates our work, which combines the strengths of TLBO and HHO while explicitly considering AoI as a primary optimization objective.

\section{System Model and Problem Statement}

\subsection{Network Architecture}
We consider a typical three-layer MEC network architecture, as shown in Fig.~\ref{fig1}. This architecture reflects the hierarchical distribution of computing resources in actual deployment, consistent with the deployment patterns of commercial MEC platforms (such as AWS Wavelength and Azure Edge Zones).

\begin{figure}[!t]
    \centering
    \centerline{\includegraphics[width=0.9\linewidth]{mec_image.png}}
    \caption{device-edge-cloud network architecture.}
    \label{fig1}
\end{figure}

The system contains the following components:
\begin{itemize}
    \item Device layer: $\mathcal{D} = \{1, 2, ..., D\}$ represents $D = 20$ mobile devices, including low-end IoT devices (40\%, 0.2-0.5 GHz), mid-range mobile devices (40\%, 0.8-1.5 GHz), and high-end devices (20\%, 2.0-3.0 GHz)
    \item Edge layer: $\mathcal{E} = \{1, 2, ..., E\}$ represents $E = 4$ edge servers (4.0-6.0 GHz)
    \item Cloud layer: $\mathcal{C} = \{1, 2, ..., C\}$ represents $C = 2$ cloud servers (8.0-12.0 GHz)
    \item Task set: $\mathcal{T} = \{1, 2, ..., T\}$ represents $T = 40$ tasks to be processed, including compute-intensive (30\%), data-intensive (25\%), real-time sensitive (25\%), and lightweight (20\%)
\end{itemize}

\subsection{Computation Model}
\subsubsection{Task Characteristics}
Each task $i \in \mathcal{T}$ is characterized by a five-tuple:
\begin{equation}
\tau_i = (D_i, C_i, T_i^{\max}, \lambda_i, AoI_i^{\max})
\end{equation}
where:
\begin{itemize}
    \item $D_i$ is the input data size (MB), $D_i^{result}$ is the output data size
    \item $C_i$ is the computational complexity (CPU cycles/bit)
    \item $T_i^{\max}$ is the maximum tolerable delay (s)
    \item $\lambda_i$ is the task update interval (s)
    \item $AoI_i^{\max}$ is the maximum acceptable age of information (s)
\end{itemize}

\subsubsection{Delay Model}
The total task processing delay includes transmission delay, waiting delay, and computation delay:
\begin{equation}
T_i = T_i^{trans} + T_i^{wait} + T_i^{comp}
\end{equation}

Local execution delay only includes computation time:
\begin{equation}
T_{i,d(i)}^{local} = \frac{D_i \cdot C_i}{f_{i,d(i)}}
\end{equation}
where $f_{i,d(i)}$ is the CPU frequency allocated to task $i$.

The delay for offloading a task to an edge server:
\begin{equation}
T_{i,d(i),e(i)}^{edge} = \frac{D_i}{R_{d(i),e(i)}} + T_{queue}^{edge} + \frac{D_i \cdot C_i}{f_{i,e(i)}} + \frac{D_i^{result}}{R_{e(i),d(i)}}
\end{equation}
where $R_{j,k}$ is the transmission rate from node $j$ to $k$, calculated using Shannon's formula:
\begin{equation}
R_{j,k} = B_{j,k} \log_2\left(1 + \frac{P_j^{trans} \cdot h_{j,k}}{N_0 \cdot B_{j,k}}\right)
\end{equation}

\textbf{Queuing delay analysis}: Considering edge servers adopt the M/M/1 queuing model, the average waiting time is:
\begin{equation}
T_{queue}^{edge} = \frac{\rho}{f_e \cdot (1-\rho)}
\end{equation}
where $\rho = \lambda_e / \mu_e$ is the server utilization rate, $\lambda_e$ is the task arrival rate, and $\mu_e$ is the service rate.

\subsubsection{Energy Model}
Computation energy consumption adopts a dynamic power model:
\begin{equation}
E_{i,j}^{comp} = \kappa_j \cdot (f_{i,j})^{\alpha} \cdot D_i \cdot C_i
\end{equation}
where $\kappa_j \in [1.0, 3.5] \times 10^{-27}$ and $\alpha \in [2.2, 2.8]$ are set according to device type.

Transmission energy consumption: $E_{i,j,k}^{trans} = P_j^{trans} \cdot (D_i + D_i^{result})/R_{j,k}$, where $P_j^{trans} \in [0.1, 0.5]$ W.

The total energy consumption of mobile devices:
\begin{equation}
    \begin{split}
        E_i = x_{i}^{local} \cdot E_{i,d(i)}^{comp} + x_{i}^{edge} \cdot E_{i,d(i),e(i)}^{trans} \\ + x_{i}^{cloud} \cdot (E_{i,d(i),e(i)}^{trans} + E_{i,e(i),c(i)}^{trans})
    \end{split}
\end{equation}

\subsubsection{AoI Model}
To describe the freshness of information, we define the Age of Information (AoI) metric. For task $i$, its AoI at time $t$ is defined as:
\begin{equation}
AoI_i(t) = t - U_i(t)
\end{equation}
where $U_i(t)$ is the time when task $i$ last successfully completed an update.

Considering task processing delay and update interval, the average AoI of task $i$ is:
\begin{equation}
\overline{AoI}_i = \lambda_i + \frac{T_i}{2} + \frac{T_i^2}{2\lambda_i}
\end{equation}

This formula considers three parts: update interval $\lambda_i$, average waiting time $T_i/2$, and additional aging due to processing delay $T_i^2/(2\lambda_i)$.

\subsection{Optimization Problem Modeling}

\subsubsection{Decision Variables}
We define binary decision variables $\mathbf{x} = \{x_i^{local}, x_i^{edge}, x_i^{cloud}\}_{i \in \mathcal{T}}$, where:
\begin{itemize}
    \item $x_{i}^{local}$ indicates task $i$ is executed locally on the device
    \item $x_{i}^{edge}$ indicates task $i$ is offloaded to an edge server
    \item $x_{i}^{cloud}$ indicates task $i$ is offloaded to a cloud server
\end{itemize}
Continuous decision variables $\mathbf{f} = \{f_{i,j}\}_{i \in \mathcal{T}, j \in \mathcal{D} \cup \mathcal{E} \cup \mathcal{C}}$ represent the allocated CPU frequency.

\subsubsection{Objective Function}
Based on practical application requirements and recent research trends, we use the weighted sum method to transform the tri-objective into a single objective:
\begin{equation}
    \begin{split}
        \min_{\mathbf{x}, \mathbf{f}} F = \sum_{i \in \mathcal{T}} \bigg( w^E \cdot \frac{E_i}{E^{norm}} + w^T \cdot \frac{T_i}{T^{norm}}
        \\+ w^{AoI} \cdot \frac{\overline{AoI}_i}{AoI^{norm}} \bigg)
    \end{split}
\end{equation}
where:
\begin{itemize}
    \item Weight settings: $w^E = 0.25$, $w^T = 0.35$, $w^{AoI} = 0.40$ (prioritizing information timeliness)
    \item Normalization factors: $E^{norm} = \max_{i \in \mathcal{T}} E_i^{local}$, $T^{norm} = \max_{i \in \mathcal{T}} T_i^{cloud}$, $AoI^{norm} = \max_{i \in \mathcal{T}} AoI_i^{\max}$
\end{itemize}

\subsubsection{Constraints}
The main constraints include:
(1) Task allocation uniqueness: $x_{i}^{local} + x_{i}^{edge} + x_{i}^{cloud} = 1, \forall i$;
(2) Resource capacity: $\sum_{i} x_{i}^{j} \cdot f_{i,j} \leq F_{j}^{\max}$;
(3) Delay constraint: $T_i \leq T_i^{\max}$;
(4) AoI constraint: $\overline{AoI}_i \leq AoI_i^{\max}$;
(5) QoS constraint: 95\% of tasks meet delay requirements.

\subsubsection{Problem Complexity Analysis}
\textbf{Theorem 1}: The MEC task offloading and resource allocation problem considering AoI is NP-hard.

\textit{Proof}: Consider a special instance: only one edge server ($E = 1$), all tasks have the same data size $D$ and computational complexity $C$, and transmission delay and AoI constraints are ignored. In this case, the problem is equivalent to:
\begin{equation}
    \begin{split}
        \min_{S \subseteq \mathcal{T}} \sum_{i \in S} p_i, \quad
        \text{s.t.} \quad p_i = \frac{DC}{f_i}, \forall i \in S, \\
        \quad
        \sum_{i \in S} f_i \leq F^{\max}, \quad
        f_i \geq 0, \forall i \in S
    \end{split}
\end{equation}
where $S$ is the set of tasks selected for edge execution, and $F^{\max}$ is the total available CPU cycles of the server.

Define the "weight" $w_i = f_i$ and "value" $v_i$ for each task $i$ as:
\begin{equation}
v_i = \frac{1}{p_i} = \frac{f_i}{DC}
\end{equation}
Then equation (14) transforms into the classic Knapsack problem of selecting items in a knapsack with capacity $F^{\max}$ to maximize total value, which has been proven to be NP-hard.

Furthermore, the complete model contains both binary variables $\mathbf{x}$ and continuous variables $\mathbf{f}$, and the objective function includes nonlinear terms (such as $f_i^{\alpha}$), while also needing to consider the temporal dependencies of AoI constraints. Therefore, it belongs to Mixed Integer Nonlinear Programming (MINLP), making the solution even more difficult. 

\section{TLBO-HHO Fusion Algorithm}

The proposed TLBO-HHO fusion algorithm addresses the multi-objective optimization challenge in MEC through an innovative two-phase integrated architecture. The core design philosophy leverages the complementary strengths of both algorithms: TLBO's stable population evolution through teacher-student interactions and HHO's dynamic search capabilities inspired by hawk hunting behavior. This synergistic combination enables effective balance between global exploration and local exploitation while incorporating AoI-aware mechanisms throughout the optimization process.

To represent the complex decision space, we encode each solution $S_k$ as a mixed vector containing three components: $S_k = [\mathbf{x}_k, \mathbf{f}_k, \mathbf{a}_k]$, where $x_i^k \in \{0, 1, 2\}$ indicates the task execution location (local, edge, or cloud), $f_i^k$ represents the allocated CPU frequency, and $a_i^k$ denotes the server assignment. This encoding naturally captures both discrete offloading decisions and continuous resource allocation variables. The population initialization employs an intelligent strategy combining 30\% heuristic solutions with 70\% random solutions, as detailed in Algorithm~\ref{alg:init}. The heuristic component prioritizes tasks based on their computational requirements and AoI constraints, while the random component ensures sufficient diversity for exploration.

\begin{algorithm}[!t]
  \small
  \caption{Intelligent Initialization Strategy}
  \label{alg:init}
  \begin{algorithmic}[1]
    \REQUIRE Population size $P$
    \ENSURE Initialized population $\mathcal{P}$
    \STATE $\mathcal{P} \leftarrow \emptyset$
    \STATE $n_{heur} \leftarrow \lfloor 0.3 \times P \rfloor$ // Heuristic part
    \FOR{$i = 1$ \TO $n_{heur}$}
        \STATE $s \leftarrow \textsc{HeuristicInit}()$
        \STATE $\mathcal{P} \leftarrow \mathcal{P} \cup \{s\}$
    \ENDFOR
    \STATE $n_{rand} \leftarrow P - n_{heur}$ // Random part
    \FOR{$i = 1$ \TO $n_{rand}$}
        \STATE $s \leftarrow \textsc{RandomInit}()$
        \STATE $\mathcal{P} \leftarrow \mathcal{P} \cup \{s\}$
    \ENDFOR
    \RETURN $\mathcal{P}$
  \end{algorithmic}
\end{algorithm}

The algorithm enhances both TLBO and HHO mechanisms to better suit the MEC optimization context. For TLBO, we introduce an adaptive teaching factor $T_F = 1 + rand() \cdot (1 + e^{-\sigma_{pop}/\sigma_{max}})$ that adjusts based on population diversity, preventing premature convergence in the teacher phase. The learning phase incorporates elite individual influence to accelerate convergence toward promising regions. For HHO, we design a nonlinear escape energy function $E = 2E_0 \cdot \cos(\pi t/2T_{max}) \cdot (1 + 0.1 \sin(4\pi t/T_{max}))$ that creates oscillating exploration-exploitation patterns, enhancing the algorithm's ability to escape local optima.

The key innovation lies in the two-phase integration mechanism, which dynamically switches between TLBO and HHO based on an adaptive probability function: $p_{HHO} = p_{base} \cdot \alpha_{conv} \cdot \beta_{div}$, where $p_{base} = 0.7 - 0.6(t/T_{max})$ ensures a gradual transition from HHO-dominated exploration (70\% initially) to TLBO-dominated exploitation (90\% finally). The convergence factor $\alpha_{conv}$ and diversity factor $\beta_{div}$ further fine-tune this transition based on real-time population characteristics. This adaptive mechanism ensures robust global search in early iterations while achieving refined local optimization in later stages.

Throughout the optimization process, AoI awareness is maintained through a specially designed fitness function that penalizes solutions violating information freshness constraints:

\begin{equation}
    \begin{split}
        Fitness(X) = w^E \cdot E(X) + w^T \cdot D(X) \\ 
        + w^{AoI} \cdot AoI(X) + Penalty(X)
    \end{split}
\end{equation}

The penalty term exponentially increases for solutions exceeding AoI thresholds, effectively guiding the search toward timely task execution strategies. Algorithm~\ref{alg:tlbo_hho_aoi} presents the complete two-phase integration procedure, where each iteration adaptively selects between TLBO teaching-learning operations and HHO exploration-exploitation strategies based on current search progress and population state.

\begin{algorithm}[!t]
  \small
  \caption{TLBO-HHO Two-Phase Integration Algorithm (with AoI Optimization)}
  \label{alg:tlbo_hho_aoi}
  \begin{algorithmic}[1]
    \REQUIRE Population size $N$, Maximum iterations $T_{\max}$, Weights $w^E, w^T, w^{AoI}$
    \ENSURE Optimal solution $\mathbf{X}_{best}$
    \STATE $P \leftarrow \textsc{IntelligentInit}(N)$; $ElitePool \leftarrow \emptyset$
    \STATE Evaluate all individual fitness (including AoI)
    \FOR{$t = 1$ \TO $T_{\max}$}
        \STATE $p_{HHO} \leftarrow \textsc{AdaptiveProb}(t, P)$
        \STATE $E \leftarrow \textsc{EscapeEnergy}(t)$; $T_F \leftarrow \textsc{TeachingFactor}(P)$
        \FOR{\textbf{each} $\mathbf{X}_i \in P$}
            \IF{$rand() < p_{HHO}$}
                \IF{$|E| \geq 1$}
                    \STATE $\mathbf{X}_{new} \leftarrow \textsc{HHO\_Explore}(\mathbf{X}_i, \mathbf{X}_{best})$
                \ELSE
                    \STATE $\mathbf{X}_{new} \leftarrow \textsc{HHO\_Exploit}(\mathbf{X}_i, \mathbf{X}_{best}, E, J)$
                \ENDIF
            \ELSE
                \STATE $\mathbf{X}_{teacher} \leftarrow \textsc{SelectTeacher}(P, ElitePool)$
                \STATE $\mathbf{X}_{new} \leftarrow \textsc{TLBO\_Teach}(\mathbf{X}_i, \mathbf{X}_{teacher}, T_F)$
                \STATE $\mathbf{X}_{new} \leftarrow \textsc{TLBO\_Learn}(\mathbf{X}_{new}, P)$
            \ENDIF
            \STATE $\mathbf{X}_{new} \leftarrow \textsc{ConstraintRepair}(\mathbf{X}_{new})$
            \STATE $f_{new} \leftarrow \textsc{Evaluate}(\mathbf{X}_{new}) + \textsc{Penalty}(\mathbf{X}_{new})$
            \IF{$f_{new} < f(\mathbf{X}_i)$}
                \STATE $\mathbf{X}_i \leftarrow \mathbf{X}_{new}$; Update $ElitePool$
            \ENDIF
        \ENDFOR
    \ENDFOR
    \RETURN $\mathbf{X}_{best}$
  \end{algorithmic}
\end{algorithm}

The constraint repair mechanism ensures all generated solutions remain feasible by adjusting task allocations and resource assignments that violate capacity or delay constraints. The elite pool maintains the top 10\% of solutions discovered throughout the search, providing memory-based guidance for future iterations. This comprehensive design enables the TLBO-HHO algorithm to effectively navigate the complex MEC optimization landscape while maintaining strict adherence to AoI requirements.

\section{Performance Evaluation}

\subsection{Experimental Setup}
Based on actual MEC deployment scenarios and configurations of AWS Wavelength and Azure Edge Zones, the parameters are set as shown in Table~\ref{tab:system}.

\begin{table}[!t]
\small
\centering
\caption{MEC System Parameter Configuration}
\label{tab:system}
\begin{tabular}{|c|c|}
\hline
\textbf{Parameter} & \textbf{Value} \\
\hline
Number of devices/edge/cloud servers & 20/4/2 \\
Device/edge/cloud CPU frequency & 0.2-3/4-6/8-12 GHz \\
Transmission power range & 0.1-0.5 W \\
Bandwidth (device-edge/edge-cloud) & 10-50/100-1000 Mbps \\
\hline
\end{tabular}
\end{table}

\begin{table}[!t]
  \scriptsize
  \setlength{\tabcolsep}{4pt}
  \centering
  \caption{Task Type Distribution and Characteristics}
  \label{tab:task}
  \resizebox{\columnwidth}{!}{%
    \begin{tabular}{|c|c|c|c|c|}
       \hline
        \textbf{Task Type} & \textbf{Ratio(\%)} & \textbf{Avg Load(G)} & \textbf{Update Int.(s)} & \textbf{Max AoI(s)} \\
        \hline
        Compute-intensive & 30 & 1.36 & 0.949 & 2.549 \\
        Data-intensive & 25 & 3.23 & 1.214 & 4.114 \\
        Real-time sensitive & 25 & 0.14 & 0.617 & 0.690 \\
        Lightweight & 20 & 0.01 & 1.302 & 1.619 \\
        \hline
    \end{tabular}}
\end{table}

The optimal parameter configuration was determined through orthogonal experiments and grid search. Common parameters: population size 50, maximum iterations 150, 5 independent runs, weight settings $w^E = 0.25$, $w^T = 0.35$, and $w^{AoI} = 0.40$. TLBO-HHO specific parameters include: initial HHO probability $p_{max} = 0.7$, final HHO probability $p_{min} = 0.1$, elite pool ratio 10\%.

These parameter choices reflect the algorithm's design philosophy and optimization priorities. The weight configuration ($w^{AoI} > w^T > w^E$) prioritizes information freshness over latency and energy consumption, which aligns with the real-time requirements of MEC applications. The adaptive HHO probability decreasing from 0.7 to 0.1 ensures strong global exploration in early iterations (70\% HHO operations) while transitioning to refined local exploitation through TLBO in later stages. This gradual shift prevents premature convergence while maintaining solution quality. The 10\% elite pool preserves high-quality solutions across iterations, accelerating convergence without sacrificing diversity. The moderate population size of 50 balances computational efficiency with solution space coverage, while 150 iterations provide sufficient time for the two-phase mechanism to fully exploit both algorithms' strengths.

\subsection{Algorithm Convergence and Performance Analysis}
Fig.~\ref{fig:convergence} shows the convergence curve comparison of the four algorithms. TLBO-HHO drops rapidly in the initial phase (0-25 iterations), maintains stable improvement in the middle phase, and finally converges to 0.00102. In contrast, TLBO converges stably but slowly (0.00292), GA easily falls into local optima (0.00757), and GWO performs poorly (0.00367).

\begin{figure}[!t]
    \centering
    \centerline{\includegraphics[width=0.9\linewidth]{convergence_comparison.png}}
    \caption{Algorithm convergence curve comparison}
    \label{fig:convergence}
\end{figure}

Fig.~\ref{fig:metrics} shows the performance metrics comparison of the four algorithms. TLBO-HHO achieves an average fitness of 0.00102, improving by 65.2\%, 85.4\%, and 74.6\% compared to TLBO, GA, and GWO respectively. In terms of energy consumption, TLBO-HHO only consumes 3.19J, significantly lower than GA (31.14J). Although GA achieves the best delay (0.577s), its energy consumption is too high; GWO has the fewest AoI violations (4.6 times), but its overall performance is inferior to TLBO-HHO.

\begin{figure}[!t]
    \centering
    \centerline{\includegraphics[width=0.9\linewidth]{performance_metrics_detail.png}}
    \caption{Algorithm performance metrics comparison}
    \label{fig:metrics}
\end{figure}

\subsection{Task Allocation Strategy and AoI Performance}
Fig.~\ref{fig:allocation} shows the task allocation location comparison. TLBO-HHO adopts a highly localized strategy (97.5\% local execution), effectively reducing transmission delay and energy consumption. In contrast, GA overuses edge resources (22.5\%), and GWO has the highest cloud utilization (12.5\%).

\begin{figure}[!t]
    \centering
    \centerline{\includegraphics[width=0.9\linewidth]{task_allocation_comparison.png}}
    \caption{Algorithm task allocation comparison}
    \label{fig:allocation}
\end{figure}

Fig.~\ref{fig:aoi_analysis} shows the AoI performance for different task types. TLBO-HHO performs best for real-time sensitive tasks, with 95\% of AoI within 1 second. Through prioritized local execution and reasonable resource allocation, the AoI of all task types remains within acceptable ranges.

\begin{figure}[!t]
    \centering
    \centerline{\includegraphics[width=0.9\linewidth]{task_type_characteristics.png}}
    \caption{Task type AoI performance analysis}
    \label{fig:aoi_analysis}
\end{figure}

\section{Conclusion}

This paper proposes a two-phase integrated TLBO-HHO algorithm for MEC task offloading optimization considering Age of Information. The algorithm achieves performance breakthrough through innovative fusion mechanisms, using TLBO's teaching phase to guide global exploration and HHO's hunting strategy to provide local exploitation. Experimental results demonstrate that TLBO-HHO improves fitness values by 65.2\%, 85.4\%, and 74.6\% compared to traditional algorithms, while maintaining excellent AoI performance for real-time tasks.

The algorithm's intelligent task allocation strategies adaptively adjust execution locations according to task types, achieving 97.5\% local execution rate to minimize transmission overhead. For real-time sensitive tasks, 95\% of AoI is controlled within 1 second through prioritized resource allocation. This differentiated processing strategy meets practical MEC application requirements.

Future research directions include: (1) extending to distributed implementation to reduce communication overhead in large-scale scenarios; (2) considering user mobility and dynamic task arrivals for online optimization; (3) integrating federated learning for privacy-preserving task offloading; (4) deploying on real MEC testbeds for practical validation.

\section*{Acknowledgment}

This work was supported by the Beijing Natural Science Foundation under Grant No. L222044.

\begin{thebibliography}{00}
\bibitem{b1} Y. Mao, C. You, J. Zhang, K. Huang, and K. B. Letaief, ``A survey on mobile edge computing: The communication perspective,'' \textit{IEEE Commun. Surveys Tuts.}, vol. 19, no. 4, pp. 2322--2358, 4th Quart., 2017.
\bibitem{b2} S. Sardellitti, G. Scutari, and S. Barbarossa, ``Joint optimization of radio and computational resources for multicell mobile-edge computing,'' \textit{IEEE Trans. Signal Inf. Process. Netw.}, vol. 1, no. 2, pp. 89--103, Jun. 2015.
\bibitem{b3} L. Huang, S. Bi, and Y. J. A. Zhang, ``Deep reinforcement learning for online computation offloading in wireless powered mobile-edge computing networks,'' \textit{IEEE Trans. Mobile Comput.}, vol. 19, no. 11, pp. 2581--2593, Nov. 2020.
\bibitem{b4} S. Guo, B. Xiao, Y. Yang, and Y. Yang, ``Energy-efficient dynamic offloading and resource scheduling in mobile cloud computing,'' in \textit{Proc. IEEE INFOCOM}, San Francisco, CA, USA, Apr. 2016, pp. 1--9.
\bibitem{b5} R. V. Rao, V. J. Savsani, and D. P. Vakharia, ``Teaching--learning-based optimization: A novel method for constrained mechanical design optimization problems,'' \textit{Comput.-Aided Des.}, vol. 43, no. 3, pp. 303--315, Mar. 2011.
\bibitem{b6} A. A. Heidari, S. Mirjalili, H. Faris, I. Aljarah, M. Mafarja, and H. Chen, ``Harris hawks optimization: Algorithm and applications,'' \textit{Future Gener. Comput. Syst.}, vol. 97, pp. 849--872, Aug. 2019.
\bibitem{b7} H. Chen, A. A. Heidari, H. Chen, M. Wang, Z. Pan, and A. H. Gandomi, ``Multi-population differential evolution-assisted Harris hawks optimization: Framework and case studies,'' \textit{Future Gener. Comput. Syst.}, vol. 111, pp. 175--198, Oct. 2020.
\bibitem{b8} B. Tu, F. Wang, Y. Huo, and X. Wang, ``A hybrid algorithm of grey wolf optimizer and harris hawks optimization for solving global optimization problems with improved convergence performance,'' \textit{Sci. Rep.}, vol. 13, no. 1, p. 22909, Dec. 2023.
\bibitem{b9} C. Wang, F. R. Yu, C. Liang, Q. Chen, and L. Tang, ``Joint computation offloading and interference management in wireless cellular networks with mobile edge computing,'' \textit{IEEE Trans. Veh. Technol.}, vol. 66, no. 8, pp. 7432--7445, Aug. 2017.
\bibitem{b10} J. Liu, Y. Mao, J. Zhang, and K. B. Letaief, ``Delay-optimal computation task scheduling for mobile-edge computing systems,'' in \textit{Proc. IEEE Int. Symp. Inf. Theory (ISIT)}, Barcelona, Spain, Jul. 2016, pp. 1451--1455.
\bibitem{b11} L. Wang, K. Wang, C. Pan, W. Xu, N. Aslam, and L. Hanzo, ``Deep reinforcement learning based dynamic trajectory control for UAV-assisted mobile edge computing,'' \textit{IEEE Trans. Mobile Comput.}, vol. 21, no. 10, pp. 3536--3550, Oct. 2022.
\bibitem{b12} Z. Li and Q. Zhu, ``Genetic algorithm-based optimization of offloading and resource allocation in mobile-edge computing,'' \textit{Information}, vol. 11, no. 2, p. 83, Feb. 2020.
\bibitem{b13} L. Zang, X. Zhang, and B. Guo, ``Federated deep reinforcement learning for online task offloading and resource allocation in WPC-MEC networks,'' \textit{IEEE Access}, vol. 10, pp. 9856--9867, 2022.
\bibitem{b14} A. Ali, S. A. A. Shah, T. Al Shloul, A. Khattak, A. Ahmad, and A. Ali, ``Multiobjective Harris Hawks optimization-based task scheduling in cloud-fog computing,'' \textit{IEEE Internet Things J.}, vol. 11, no. 13, pp. 24334--24352, Jul. 2024.
\bibitem{b15} M. A. Abd-Elmagid, N. Pappas, and H. S. Dhillon, ``On the role of age of information in the Internet of Things,'' \textit{IEEE Commun. Mag.}, vol. 57, no. 12, pp. 72--77, Dec. 2019.
\bibitem{b16} S. Kaul, R. Yates, and M. Gruteser, ``Real-time status: How often should one update?'' in \textit{Proc. IEEE INFOCOM}, Orlando, FL, USA, Mar. 2012, pp. 2731--2735.
\bibitem{b17} L. Liu, J. Feng, X. Mu, Q. Pei, D. Lan, and M. Xiao, ``Asynchronous deep reinforcement learning for collaborative task computing and on-demand resource allocation in vehicular edge computing,'' \textit{IEEE Trans. Intell. Transp. Syst.}, vol. 24, no. 12, pp. 15513--15526, Dec. 2023.
\bibitem{b18} Z. Zhou, P. Liu, J. Feng, Y. Zhang, S. Mumtaz, and J. Rodriguez, ``Computation resource allocation and task assignment optimization in vehicular fog computing: A contract-matching approach,'' \textit{IEEE Trans. Veh. Technol.}, vol. 68, no. 4, pp. 3113--3125, Apr. 2019.
\end{thebibliography}

\end{document}